% !TEX root = ../report.tex
Схема линейного поиска представлена на рисунке~\ref{fig:linear}.

\begin{figure}[H]
\centering
\begin{tikzpicture}[flow]
\node[term] (start) at (0,0.0) {Начало: $a,N,x$};
\node[proc] (init) at (0,-1.6) {$i\gets0$};
\node[decision] (loop) at (0,-3.9) {$i<N$};
\node[decision] (equal) at (0,-6.7) {$a_i=x$};
\node[proc] (inc) at (0,-9.0) {$i\gets i+1$};
\node[term] (miss) at (5,-3.9) {Возврат $-1$};
\node[term] (found) at (5,-6.7) {Возврат $i$};
\draw[arrow] (start)--(init);
\draw[arrow] (init)--(loop);
\draw[arrow] (loop)--node[left]{Да}(equal);
\draw[arrow] (loop)--node[above]{Нет}(miss);
\draw[arrow] (equal)--node[above]{Да}(found);
\draw[arrow] (equal)--node[left]{Нет}(inc);
\draw[arrow] (inc.west)--(-3,-9.0)--(-3,-2.5)--(0,-2.5)--(loop.north);
\end{tikzpicture}
\caption{Схема линейного поиска}\label{fig:linear}
\end{figure}
